\begin{textsample}{NEG Dim 2 – al – Score 1.00 – al/al\_ef\_68.txt}  \label{ex:f2_neg_038}
3.4 Modelagem Matemática Modelagem Matemática está relacionada com a resolução de problemas e com os procedimentos e tem como objetivo obter um modelo matemático que seja a solução do problema inicial.
Na Educação Básica, a Modelagem é vista como uma alternativa pedagógica na qual é utilizada uma situação problema real ou da própria Matemática.
Alguns aspectos são importantes no desenvolvimento dos trabalhos com modelagem : investigação autônoma ( trabalho em grupo ), ciclo de modelagem e temas do mundo real.
A ideia é fazer modelagem para aprender matemática, ou seja, os Modelos Matemáticos precisam ser significativos com situações-problema úteis e possíveis de serem resolvidas e discutidas. “ Modelagem Matemática é um processo dinâmico utilizado para a obtenção e validação de modelos matemáticos.
É uma forma de abstração e generalização com a finalidade de previsão de tendências.
A modelagem consiste, essencialmente, na arte de transformar situações da realidade em problemas matemáticos cujas soluções devem ser interpretadas na linguagem usual ” ( Bassanezi, 2004, p.24 ).
Segundo Rigonatto, a modelagem matemática é a criação de um modelo matemático ( um padrão ou fórmula matemática ) para explicação ou compreensão de um fenômeno natural.
Esse fenômeno pode ser de qualquer área do conhecimento.
Através do uso da modelagem matemática na sala de aula podemos trabalhar a interdisciplinaridade, a transversalidade, mostrando ao estudante como a matemática pode ser útil em sua vida fora do ambiente escolar e como ela interage com as demais áreas do conhecimento.
O estudante passa a perceber a importância da matemática para a compreensão de fenômenos naturais, como é possível “ prever ” alguns acontecimentos utilizando fórmulas e modelos e isso acaba despertando seu interesse pela ciência.
A introdução da modelagem matemática pode ser feita através da resolução de problemas, trazendo para dentro de sala a realidade do estudante.
As diversas situações-problema farão com que a capacidade de interpretação melhore, o estudante assuma uma posição crítica ao tentar resolvê -las e consiga analisar que pode haver mais de uma solução e que há vários caminhos para chegar até elas.

% matched lemmas: 
\end{textsample}
